%% LyX 2.4.4 created this file.  For more info, see https://www.lyx.org/.
%% Do not edit unless you really know what you are doing.
\documentclass[english]{article}
\usepackage[T1]{fontenc}
\usepackage[latin9]{inputenc}
\usepackage{enumitem}
\usepackage{amsmath}
\usepackage{amsthm}
\usepackage{amssymb}
\usepackage{geometry}
\geometry{verbose,tmargin=1in,bmargin=1in,lmargin=1in,rmargin=1in}
\PassOptionsToPackage{normalem}{ulem}
\usepackage{ulem}

\makeatletter
%%%%%%%%%%%%%%%%%%%%%%%%%%%%%% Textclass specific LaTeX commands.
\numberwithin{equation}{section}
\numberwithin{figure}{section}
\newlength{\lyxlabelwidth}      % auxiliary length 

%%%%%%%%%%%%%%%%%%%%%%%%%%%%%% User specified LaTeX commands.
\usepackage{fancyhdr}
\usepackage{lastpage}
\pagestyle{fancy}
\usepackage{enumitem}
\usepackage{ifthen}
\everymath=\expandafter{\the\everymath\displaystyle}
%\DeclareMathSizes{10}{10}{10}{10}
\usepackage{multicol}

\usepackage{tikz}
\usetikzlibrary{patterns}
\usetikzlibrary{plotmarks}
\usepackage{pgfplots}
\pgfplotsset{compat=newest}

\usepackage{calc}
\usepackage[nomessages]{fp}

\usepackage{datetime}


%%%%%commands
\newcommand{\mb}[1]{\mathbb{#1}}
\newcommand{\funct}[3]{$#1\colon#2\rightarrow#3$}
% Added by lyx2lyx
\setlength{\parskip}{\smallskipamount}
\setlength{\parindent}{0pt}

\makeatother

\usepackage{babel}
\begin{document}
\renewcommand{\author}{Dr.\,James Ashe}

\newcommand{\classNum}{439}

\newcommand{\classSec}{A}

\newcommand{\nameType}{Name: \hspace{-4cm}}

\newcommand{\examType}{Midterm}

\newcommand{\Date}{\formatdate{4}{3}{2014}} %%\AdvanceDate[12] \today

%\newdate{my_date}{\day}{\month}{\year}%   
%\AdvanceDate[-5] 
%\displaydate{my_date}

%%First page's header%%%%%%%%%%%%%%%%%%%%%%
\fancypagestyle{firststyle} %%header settings for first pagr
{
	\fancyhead[L]{MTH \classNum{}(\classSec{}) \author{}\\
	\textbf{\examType{}} \Date{}}
	\fancyhead[C]{\textbf{\nameType}}
	\setlength{\headsep}{14pt}
}
%%%%%%%%%%%%%%%%%%%%%%%%%%%%%%%%%%%%%%%%%%

%%Next page's header%%%%%%%%%%%%%%%%%%%%%%%
%\setlist{nolistsep} %eliminates space between list items
\lhead{MTH \classNum{}(\classSec{})} 
\chead{\textbf{\nameType}} 
\cfoot{\thepage{}~of~\pageref{LastPage}} 
\setlength{\headsep}{5pt} 
%%%%%%%%%%%%%%%%%%%%%%%%%%%%%%%%%%%%%%%%%%%

%%Various settings; see the preamble for other settings%%%%%
%%\setenumerate[0]{leftmargin=12pt,itemindent=0pt,labelwidth=0pt}
\renewcommand{\labelenumi}{\textbf{\arabic{enumi}.}}
\renewcommand{\labelenumii}{\textbf{\alph{enumii}.}}
\thispagestyle{firststyle}
%%%%%%%%%%%%%%%%%%%%%%%%%%%%%%%%%%%%%%%%%%
\newcommand{\xspace}{\vspace{.85cm}}

\emph{Unless stated otherwise, all rings are assumed to have the normal
addition and multiplication operations.}

\subsection*{Part 1}
\begin{enumerate}[resume]
\item {[}3 pts each{]} Provide the definition for each of the following:

\begin{enumerate}
\item Integral domain\xspace
\item Onto function\xspace
\item Ring isomorphism\xspace
\item Ideal\xspace\vfill
\end{enumerate}
\item {[}3 pts each{]} Provide an example for each of the following (do
not use an example from problems \ref{enu:first}-\ref{enu:last}):

\begin{enumerate}
\item A ring without identity\xspace
\item A non-commutative ring \xspace
\item A field \xspace
\item A function which is not 1-1\xspace
\item An onto function\xspace
\item Two rings which are isomorphic\xspace
\item Two rings which are not isomorphic\xspace
\item An ideal\xspace\vfill
\end{enumerate}
\end{enumerate}
\emph{\newpage}

\subsection*{Part 2}

\emph{{[}9 pts each{]} Choose the indicated number of problems from
each section.}

\subsubsection*{Rings: choose any three.}
\begin{enumerate}
\item \label{enu:first}Complete the following:

\begin{enumerate}[resume]
\item Give an example of a subset of the integers which \uline{is not} a
ring (under the usual addition and multiplication).
\item Give an example of a subset of the integers which \uline{is} a ring.
\end{enumerate}
\end{enumerate}
\begin{enumerate}[resume]
\item Prove that $\mathbb{Q}$ is subring of $\mathbb{R}$.
\item Complete the following:

\begin{enumerate}[resume]
\item Prove that the set of odd integers is not a subring of $\mathbb{Z}$.
\item Prove that in general the union of any two rings is not a ring.
\end{enumerate}
\item Prove that the intersection of any two rings is also a ring.
\item Let $R$ be a ring and $a,b\in R$. Prove that $(-a)(-b)=ab$. Hint:
$(a+(-a))(-b)=0$.
\end{enumerate}
\vfill

\subsubsection*{Integral Domains and Fields: choose any two.}
\begin{enumerate}[resume]
\item Prove that $\mathbb{Z}_{n}$ is an integral domain if $n$ is not
prime.
\item Prove that $\mathbb{Z}_{3}$ is a field. 
\item Let $R$ be an integral domain and suppose that $ax-bx=0$. Prove
that if $a\neq b$ then $x=0$.
\item Let $R$ be an integral domain and let $a,b,c\in R$ where $c\neq0$.
If $ac=bc$ then $a=b$.
\item Let $R$ be a field. Prove that $R$ is an integral domain.
\end{enumerate}
\vfill

\subsubsection*{Homomorphisms, isomorphisms, and ideals: choose any three.}
\begin{enumerate}[resume]
\item Prove that $f\colon\mathbb{R}\rightarrow\mathbb{R}$ where $f(x)=e^{x}$
is one-to-one but is not onto.
\item Let \funct{f}{R}{S} and \funct{g}{S}{T} for rings $R$, $S$, and
$T$. Prove that \funct{f\circ g}{R}{T} is a homomorphism.
\item Let $f\colon\mathbb{Z}_{4}\rightarrow\mathbb{Z}_{4}$. 

\begin{enumerate}[resume]
\item If $f(x)=2x$ prove that $f$ is not a homomorphism.
\item If $f(x)=3x$ prove that $f$ is a homomorphism.
\end{enumerate}
\item Prove that $f\colon\sqrt{2}\mathbb{Z}\rightarrow\sqrt{2}\mathbb{Z}$
where $f(a+b\sqrt{2})=a-b\sqrt{2}$ is an isomorphism ($\sqrt{2}\mathbb{Z}=\{a+b\sqrt{2}\colon a,b\in\mathbb{Z}\}$).
\item \label{enu:last}Show that any ring isomorphism \funct{f}{\mb{Z}}{\mb{Z}}
must be the identity map.
\item Let $A$ and $B$ be ideals of $R$. Prove that $A\cap B$ is an ideal
of $R$.
\item Prove that $\mathbb{R}\times\mathbb{R}=\{(a,b)\colon a,b\in\mathbb{R}\}$
is isomorphic to $\mathbb{C}=\{a+b\sqrt{-1}\colon a,b\in\mathbb{R}\}$
if multiplication in $\mathbb{R}\times\mathbb{R}$ is defined as $(a,b)(c,d)=(ac-bd,ad+bc)$
and addition is defined as $(a,b)+(c,d)=(a+c,b+d)$.
\end{enumerate}
\vfill
\end{document}
